\chapter{Ammonia Levels}

\section*{What Is This Test?}
The ammonia (NH\textsubscript{3}) blood test measures the concentration of ammonia in the blood. Ammonia is a potent neurotoxin produced primarily in the gastrointestinal tract by bacterial degradation of proteins and amino acids, and to a lesser extent from skeletal muscle catabolism and renal production. Under normal conditions, ammonia is rapidly transported via the portal vein to the liver, where it is converted to non-toxic urea through the urea cycle (ornithine cycle) in periportal hepatocytes: 2 NH\textsubscript{3} + CO\textsubscript{2} $\rightarrow$ Urea + H\textsubscript{2}O (net reaction). A small portion of ammonia is also used by perivenous hepatocytes for glutamine synthesis. Urea is then excreted by the kidneys. In liver failure or portosystemic shunting, ammonia bypasses hepatic metabolism, enters the systemic circulation, crosses the blood-brain barrier, and is metabolized by astrocytes to glutamine (via glutamine synthetase). Glutamine accumulation causes astrocyte swelling and cerebral edema, manifesting clinically as hepatic encephalopathy. The relationship between serum ammonia levels and hepatic encephalopathy severity is imperfect: ammonia levels correlate poorly with the degree of encephalopathy, and a normal ammonia level does not exclude hepatic encephalopathy. However, markedly elevated ammonia >150 $\mu$mol/L in acute liver failure indicates poor prognosis and is a criterion for liver transplantation evaluation. In neonates, hyperammonemia >150 $\mu$mol/L is a metabolic emergency requiring immediate investigation for urea cycle disorders. Ammonia is measured in plasma, most commonly from venous blood collected in an EDTA (purple-top) or heparin tube, placed on ice immediately, and analyzed within 15--30 minutes of collection.

\section*{Alternative Names}
NH3; Serum Ammonia; Plasma Ammonia; Blood Ammonia; Hyperammonemia Screen; Ammonia, Plasma; NH4 (ammonium)

\section*{Principle}
Plasma ammonia is measured by an enzymatic method using glutamate dehydrogenase (GLDH): Alpha-ketoglutarate + NH\textsubscript{4}\textsuperscript{+} + NADPH $\rightarrow$ Glutamate + NADP\textsuperscript{+} + H\textsubscript{2}O. The rate of decrease in absorbance at 340 nm (as NADPH is oxidized to NADP\textsuperscript{+}) is proportional to ammonia concentration. The GLDH method is specific, rapid, and is used on most automated chemistry analyzers. Some point-of-care and blood gas analyzers use an ammonium ion-selective electrode (ISE) with a gas-permeable ammonia membrane. Sample handling is the most critical aspect of the test: ammonia concentration increases significantly in vitro after blood is drawn because red blood cells produce ammonia through deamination of amino acids and the urea cycle enzyme arginase. Blood must be collected without prolonged tourniquet time, placed on ice immediately, transported to the laboratory within 15 minutes, centrifuged at 4 degrees C, and analyzed within 30 minutes of collection. Sodium fluoride tubes inhibit in vitro ammonia production to some degree and may be acceptable. Samples not adhering to strict collection and handling protocols are often rejected by the laboratory because delayed analysis produces falsely elevated results.

\section*{History}
The urea cycle was discovered by Hans Krebs and Kurt Henseleit in 1932 \cite{krebs1932untersuchungen}. The neurotoxicity of ammonia was recognized in the 19th century in studies of ``meat intoxication'' in dogs with portocaval shunts (Eck fistula, 1877). The clinical association between hyperammonemia and hepatic encephalopathy was established in the 1950s. The first ammonia assays using the Conway microdiffusion method or ion-exchange were developed in the 1930s--1950s. Da Fonseca-Wollheim developed the GLDH enzymatic method in 1973, which remains the standard today. The appreciation that ammonia levels correlate poorly with hepatic encephalopathy severity has tempered enthusiasm for routine measurement, but ammonia remains essential in neonatal metabolic emergencies and acute liver failure.

\section*{Why This Test Is Ordered}
\begin{itemize}
  \item Adjunctive evaluation of hepatic encephalopathy in patients with known or suspected liver disease and altered mental status: elevated ammonia supports the diagnosis, but a normal ammonia does NOT exclude it. The diagnosis of hepatic encephalopathy is primarily clinical
  \item Assessment of acute liver failure prognosis: arterial ammonia >150 $\mu$mol/L (approximately 255 $\mu$g/dL) at presentation predicts high risk of cerebral edema and intracranial hypertension \cite{clemmesen1999cerebral}, and is used as a criterion for urgent liver transplant evaluation in some centers (King's College criteria)
  \item Diagnosis of urea cycle disorders in neonates, infants, or children presenting with lethargy, vomiting, seizures, respiratory alkalosis, and unexplained coma: plasma ammonia is the first and most important screening test \cite{haeberle2012suggested}. Levels >150 $\mu$mol/L in neonates are a metabolic emergency requiring immediate intervention to prevent irreversible neurological damage
  \item Evaluation of suspected Reye syndrome in children with encephalopathy and hepatic dysfunction following viral illness (especially varicella or influenza) and aspirin use
  \item Detection of valproic acid-induced hyperammonemia: valproic acid inhibits carbamoyl phosphate synthetase I (CPS-I) and reduces carnitine levels, causing hyperammonemia in 30--50\% of patients \cite{wadzinski2007valproate}. Often asymptomatic but can cause encephalopathy, especially when used with other antiepileptics (especially topiramate, phenobarbital)
  \item Diagnosis of portosystemic shunting (e.g., in cirrhosis with spontaneous shunts or post-TIPS procedure): ammonia is elevated due to bypass of hepatic metabolism
  \item Monitoring of therapy for urea cycle disorders (sodium benzoate, sodium phenylacetate, glycerol phenylbutyrate, nitrogen scavengers) and hepatic encephalopathy (lactulose, rifaximin)
\end{itemize}

\section*{Primary vs Secondary Test}
Ammonia is a secondary, specialized test. It is not part of routine panels. It is used adjunctively in hepatic encephalopathy (where clinical diagnosis is primary) and is the primary screening test for urea cycle disorders in neonates and children.

\section*{How This Test Is Performed}
Blood for ammonia must be collected with meticulous attention to pre-analytical conditions: (1) The patient should be at rest and avoid fist clenching (muscle activity produces ammonia). (2) A tourniquet is applied briefly; prolonged tourniquet time with venous stasis causes local ammonia accumulation from muscle metabolism. (3) Blood is drawn into an EDTA (purple-top) or lithium heparin (green-top) tube. The tube must be filled completely. (4) The tube is placed on ice IMMEDIATELY and transported to the laboratory within 15 minutes. (5) The sample is centrifuged at 4 degrees C (refrigerated centrifuge) and the plasma is separated from cells. (6) The plasma ammonia is analyzed within 30 minutes of collection. Samples not processed within this window will have falsely elevated ammonia and may be rejected by the laboratory. Arterial ammonia is slightly more accurate than venous (peripheral tissue metabolism can elevate venous ammonia), but venous is adequate for most clinical indications. For arterial ammonia, an arterial blood gas syringe kept on ice is used.

\section*{What Preparation Is Needed}
The patient should avoid smoking for at least 2 hours before the test (cigarette smoke contains ammonia and smoking increases endogenous ammonia production). Fasting for 8 hours is recommended, as a protein meal can increase ammonia levels modestly. Inform the provider of medications: valproic acid (markedly increases ammonia), topiramate (can increase ammonia), carbamazepine, phenobarbital, and certain chemotherapeutic agents (asparaginase, which depletes asparagine and glutamine, causing secondary hyperammonemia). Recent blood transfusion within 24 hours (ammonia in stored blood increases with storage duration). Recent GI bleeding (blood protein load $\rightarrow$ increased ammonia production in gut). Recent high-intensity exercise elevates ammonia. The patient should be calm and at rest during phlebotomy. The nursing and phlebotomy staff must be specifically trained in ammonia collection protocol---this is one of the most pre-analytically demanding laboratory tests.

\section*{Contraindications and Precautions}
\begin{itemize}
  \item No absolute contraindications. Critical pre-analytical considerations (the most important aspect of this test):
  \item (1) In vitro ammonia generation in blood after collection is the primary source of error. Red blood cells contain arginase and other enzymes that decompose amino acids, producing ammonia that accumulates at a rate of approximately 0.5--1 $\mu$mol/L/hour at room temperature. Immediate cooling on ice and rapid centrifugation are essential.
  \item (2) Fist clenching and muscle activity during phlebotomy: ammonia is produced by exercising muscle (adenylate deaminase pathway) and can cause a localized increase in the drawn limb.
  \item (3) Prolonged tourniquet time >1 minute: venous stasis concentrates ammonia locally.
  \item (4) Hemolysis: red blood cell ammonia is 3--5$\times$ higher than plasma ammonia. Even slight hemolysis can significantly elevate measured ammonia.
  \item (5) Ammonia contamination from laboratory glassware or reagents---most modern clinical labs use ammonia-free water and dedicated ammonia chemistry channels to avoid contamination.
  \item (6) Delayed processing: the longer the sample sits at room temperature before centrifugation, the higher the ammonia will be. A sample sitting at room temperature for 1 hour before processing will have a falsely elevated ammonia, potentially 50--100\% higher than the true value.
  \item (7) Drugs causing hyperammonemia: valproic acid (inhibits CPS-I and reduces carnitine---can cause severe hyperammonemia and encephalopathy at therapeutic valproate levels; check ammonia in all patients on valproate with altered mental status), asparaginase chemotherapy (depletes asparagine/glutamine), 5-fluorouracil, salicylates (in Reye syndrome), and topiramate (mild hyperammonemia, synergistic with valproate).
  \item (8) Ammonia levels show poor correlation with hepatic encephalopathy severity. In patients with cirrhosis and altered mental status, the diagnosis of hepatic encephalopathy is clinical, and a normal ammonia does NOT rule it out. Do not withhold lactulose/rifaximin therapy for suspected hepatic encephalopathy while awaiting ammonia levels.
  \item (9) In acute liver failure, arterial ammonia is preferred (avoids artifact from peripheral muscle metabolism) and is a key prognostic marker.
\end{itemize}
\section*{Risks}
Minimal risk from venipuncture. For arterial ammonia (drawn from radial or femoral artery), there is a slightly higher risk of hematoma, arterial dissection, and digital ischemia (very rare with proper technique).

\section*{What Happens After the Test}
Results available within 1--2 hours. Interpretation requires consideration of the clinical context and strict attention to sample handling: (1) If the sample was not processed correctly (delayed, not on ice, tourniquet prolonged), the ammonia level may be falsely elevated and should be interpreted with extreme caution or repeated. (2) Elevated ammonia in a patient with cirrhosis and altered mental status supports hepatic encephalopathy---but clinical improvement with lactulose/rifaximin is the diagnostic endpoint. Ammonia levels are not useful for monitoring treatment response. (3) Markedly elevated ammonia >150 $\mu$mol/L (>255 $\mu$g/dL) in acute liver failure: indicates high risk of cerebral edema and intracranial hypertension; urgent hepatology consultation for transplant evaluation. (4) In neonates with lethargy, vomiting, seizures, and hyperammonemia >150 $\mu$mol/L: assume urea cycle disorder until proven otherwise. Stop all protein intake immediately. Initiate IV glucose (10\%) to prevent catabolism. Start ammonia scavengers (sodium benzoate + sodium phenylacetate, or glycerol phenylbutyrate). Arrange urgent hemodialysis if ammonia >400--500 $\mu$mol/L or not rapidly responsive to medical therapy. Collect plasma amino acids and urine orotic acid for specific diagnosis. (5) Normal ammonia with clinical encephalopathy: does not exclude hepatic encephalopathy; up to 30--40\% of patients with cirrhosis and hepatic encephalopathy have normal ammonia levels.

\section*{Understanding Results}

\subsection*{Below Normal}
\begin{itemize}
  \item Low ammonia is not clinically significant and requires no evaluation.
  \item Ammonia normally exists at very low concentrations (<35 $\mu$mol/L). Low levels simply reflect normal ammonia metabolism.
  \item In hepatic encephalopathy, ammonia may be normal (especially in chronic hepatic encephalopathy).
  \item Some antihypertensive medications (sodium nitroprusside) may transiently lower ammonia, but this is not clinically significant.
\end{itemize}

\subsection*{Above Normal}
\begin{itemize}
  \item \textbf{Elevated plasma ammonia (>35 $\mu$mol/L or >50 $\mu$mol/L depending on laboratory reference):} Must be interpreted with strict consideration of sample handling.
  \item \textbf{Hepatic encephalopathy (most common cause in adults):} Ammonia produced by gut bacteria cannot be metabolized by the diseased liver, and portosystemic shunts allow ammonia to bypass the liver entirely \cite{vilstrup2014hepatic}. Ammonia enters the brain, where astrocytes convert it to glutamine, causing osmotic swelling and cerebral edema. Ammonia levels commonly range from 40--200 $\mu$mol/L in cirrhosis. However, correlation with encephalopathy severity is poor---ammonia levels do not reliably distinguish mild (West Haven grade 1) from severe (grade 3--4) encephalopathy. Diagnosis remains clinical (disorientation, asterixis, fetor hepaticus, hyperreflexia, progression to stupor and coma). Treatment: lactulose (acidifies colon $\rightarrow$ NH\textsubscript{3} trapped as NH\textsubscript{4}\textsuperscript{+}), rifaximin (non-absorbable antibiotic reducing ammonia-producing gut bacteria).
  \item \textbf{Urea cycle disorders (most important use in children and neonates):} Inborn errors of the urea cycle: ornithine transcarbamylase (OTC) deficiency (most common, X-linked, hyperammonemia + elevated urine orotic acid), carbamoyl phosphate synthetase I (CPS-I) deficiency (low orotic acid), argininosuccinate synthetase deficiency (citrullinemia type I---massively elevated citrulline), argininosuccinate lyase deficiency (argininosuccinic aciduria), arginase deficiency (hyperargininemia, spastic diplegia). Presentation: healthy neonate becomes lethargic, vomiting, hypothermia, hyperventilation (respiratory alkalosis from central hyperventilation due to ammonia stimulation of respiratory center), progressing to seizures and coma within 24--72 hours of birth (after protein feeding begins). Ammonia >150 $\mu$mol/L emergently. Treatment: stop protein, IV D10W to prevent catabolism, ammonia scavengers, and urgent hemodialysis if NH3 >400--500 or neurological deterioration. Late-onset forms present in childhood/adolescence with episodic encephalopathy triggered by infection, high-protein meal, or valproate.
  \item \textbf{Acute liver failure:} Ammonia >150 $\mu$mol/L (approximately 255 $\mu$g/dL) on arterial measurement is associated with high risk of cerebral edema, intracranial hypertension, and brain herniation. Used as a transplant listing criterion (King's College criteria for non-acetaminophen ALF). Ammonia level is a better predictor of cerebral edema than grade of encephalopathy. Treatment: mannitol, hypertonic saline, therapeutic hypothermia (controversial), and expedited transplant evaluation.
  \item \textbf{Valproic acid-induced hyperammonemia:} 30--50\% incidence. Mechanism: valproate inhibits CPS-I, decreases hepatic N-acetylglutamate (an essential allosteric activator of CPS-I), and causes secondary carnitine deficiency. Usually asymptomatic (mild elevation, 40--80 $\mu$mol/L) but can cause severe hyperammonemic encephalopathy (especially with concurrent use of topiramate or phenobarbital). Treatment: discontinue valproate, administer L-carnitine (100 mg/kg/day IV), supportive care.
  \item \textbf{Reye syndrome:} Children with recent viral illness (varicella, influenza) who received aspirin $\rightarrow$ fulminant hepatic failure with hyperammonemia, hypoglycemia, and cerebral edema. Now rare since aspirin is avoided in children.
  \item \textbf{Other causes:} GI bleeding (protein load from blood in gut), high-protein diet, TPN with high amino acid content, ureterosigmoidostomy (urinary ammonia absorbed across colonic mucosa), severe infection/sepsis, portosystemic shunts (spontaneous or TIPS), chemotherapy (asparaginase, 5-FU), blood transfusion with older stored blood (ammonia content increases with storage), multiple myeloma (unknown mechanism), and sample handling artifact (delayed processing, not on ice, prolonged tourniquet).
\end{itemize}

\section*{Normal Results}
\begin{itemize}
  \item \textbf{Plasma ammonia (venous, adults):} 15--45 $\mu$mol/L (25--78 $\mu$g/dL) --- varies by laboratory
  \item \textbf{Plasma ammonia (arterial, adults):} 12--35 $\mu$mol/L (slightly lower than venous)
  \item \textbf{Newborns (term, day 1):} <80 $\mu$mol/L (higher in first days of life)
  \item \textbf{Neonatal hyperammonemia emergency threshold:} >150 $\mu$mol/L (measured; >255 $\mu$g/dL)
  \item \textbf{Acute liver failure poor prognosis (arterial):} >150 $\mu$mol/L
  \item \textbf{Conversion:} 1 $\mu$mol/L = 1.7 $\mu$g/dL; 1 $\mu$g/dL = 0.588 $\mu$mol/L
  \item \textbf{Critical values:} >100 $\mu$mol/L in adults; >150 $\mu$mol/L in neonates
\end{itemize}

\section*{Follow-up Tests}
\begin{itemize}
  \item \textbf{Liver function tests (ALT, AST, ALP, bilirubin, albumin, INR):} For liver disease.
  \item \textbf{Plasma amino acids (quantitative):} Essential for diagnosis of urea cycle disorders. Specific patterns: elevated citrulline (citrullinemia, argininosuccinic aciduria); elevated argininosuccinic acid (argininosuccinic aciduria); elevated arginine (arginase deficiency); low citrulline (OTC deficiency, CPS-I deficiency).
  \item \textbf{Urine orotic acid:} Elevated in OTC deficiency; low/normal in CPS-I deficiency. Key differentiating test.
  \item \textbf{Urine organic acids:} For organic acidemias that cause secondary hyperammonemia (methylmalonic acidemia, propionic acidemia).
  \item \textbf{Arterial blood gas (ABG):} Respiratory alkalosis is characteristic in urea cycle disorders and hepatic encephalopathy.
  \item \textbf{Serum glucose:} Hypoglycemia often accompanies hyperammonemia in acute liver failure and Reye syndrome.
  \item \textbf{Serum valproic acid level:} If valproate-induced hyperammonemia suspected.
  \item \textbf{Carnitine levels (total and free):} Carnitine deficiency in valproate-induced hyperammonemia.
  \item \textbf{Abdominal imaging with Doppler:} For portosystemic shunts and cirrhosis.
  \item \textbf{EEG:} In hepatic encephalopathy shows triphasic waves and diffuse slowing.
  \item \textbf{Genetic testing:} For specific urea cycle enzyme mutations.
  \item \textbf{Neuroimaging (CT/MRI brain):} For cerebral edema in acute liver failure or hyperammonemia.
\end{itemize}

\section*{References}
\bibliographystyle{unsrtnat}
\bibliography{references}

\clearpage
\section*{Regulatory Status and Authorization Date}
This chapter describes a test category, not a specific commercial product. FDA, CDSCO, EU IVDR, and MHRA approval or registration dates therefore cannot be assigned without the assay manufacturer, product name, instrument, and intended use. For laboratory-developed tests, an individual agency approval date may not exist. See the Preface for the interpretation of regulatory status.

\clearpage
